\documentclass[11pt]{article}

\usepackage[a4paper,margin=1in]{geometry}
\usepackage{amsmath,amssymb,amsthm,mathtools}
\usepackage{microtype}
\usepackage{enumitem}
\usepackage[hidelinks]{hyperref}

\newtheorem{theorem}{Theorem}[section]
\newtheorem{proposition}[theorem]{Proposition}
\newtheorem{lemma}[theorem]{Lemma}
\newtheorem{corollary}[theorem]{Corollary}
\theoremstyle{definition}
\newtheorem{definition}[theorem]{Definition}
\theoremstyle{remark}
\newtheorem{remark}[theorem]{Remark}

\newcommand{\R}{\mathbb{R}}
\newcommand{\one}{\mathbf{1}}
\newcommand{\tr}{\operatorname{Tr}}

\title{The Exact Commonality Region of Adjacent Even and Odd Cycles}
\author{}
\date{}

\begin{document}
\maketitle

\begin{abstract}
For an integer $k\ge 2$, we determine exactly the parameters for which
$(C_{2k},C_{2k+1})$ is a common pair of graphs in the sense of Behague,
Morrison, and Noel.  Let $a_k^\star\in(1/2,1)$ be the
unique solution of
\[
  \frac{(a_k^\star)^{2k-1}(a_k^\star+2k)}{2k+1}=2^{1-2k}.
\]
We prove
\[
  \pi(C_{2k},C_{2k+1})=(0,a_k^\star].
\]
The upper bound is witnessed by the balanced disjoint union of two cliques.  For
the lower bound, after passing to a finite-rank graphon operator, a rank-one
Lidskii--Karamata inequality separates the Perron eigenvalue from the remaining
spectrum.  A Hilbert--Schmidt budget shows that at most one negative eigenvalue
can contribute with the wrong sign.  Perron domination then pays for that single
dangerous mode exactly, reducing the graphon inequality to the elementary
convexity inequality $x^{2k}+(1-x)^{2k}\ge 2^{1-2k}$.
\end{abstract}

\section{Definitions and statement of the result}

A \emph{graphon} is a symmetric measurable function
$W\colon[0,1]^2\to[0,1]$.  For a finite graph $H$, its homomorphism density in
$W$ is
\[
  t(H,W)
  :=
  \int_{[0,1]^{V(H)}}
  \prod_{uv\in E(H)} W(x_u,x_v)
  \prod_{u\in V(H)}dx_u.
\]
Write $e(H)=|E(H)|$.  For background on graphons and homomorphism densities
we refer to Lov\'asz~\cite{Lovasz}.

Following Behague, Morrison, and Noel~\cite[Definition~1.1]{BMN}, for
$a,b\in(0,1)$ with
$a+b=1$, a pair $(H_1,H_2)$ is said to be \emph{$(a,b)$-common} if
\begin{equation}
  \frac{t(H_1,W)}{e(H_1)a^{e(H_1)-1}}
  +
  \frac{t(H_2,1-W)}{e(H_2)b^{e(H_2)-1}}
  \ge
  \frac{a}{e(H_1)}+\frac{b}{e(H_2)}
  \label{eq:common-definition}
\end{equation}
for every graphon $W$.  Define
\[
  \pi(H_1,H_2)
  :=
  \{a\in(0,1):(H_1,H_2)\text{ is }(a,1-a)\text{-common}\}.
\]

\begin{theorem}[Exact adjacent-cycle commonality region]
  \label{thm:main}
  Let $k\ge2$.  There is a unique number $a_k^\star\in(1/2,1)$ satisfying
  \begin{equation}
    \frac{(a_k^\star)^{2k-1}(a_k^\star+2k)}{2k+1}
    =2^{1-2k}.
    \label{eq:critical-k}
  \end{equation}
  Then
  \begin{equation}
    \boxed{\pi(C_{2k},C_{2k+1})=(0,a_k^\star].}
    \label{eq:pi-answer}
  \end{equation}
  Equivalently, for $a\in(0,1)$ and $b=1-a$, the inequality
  \begin{equation}
    \frac{t(C_{2k},W)}{2k\,a^{2k-1}}
    +
    \frac{t(C_{2k+1},1-W)}{(2k+1)b^{2k}}
    \ge
    \frac{a}{2k}+\frac{b}{2k+1}
    \label{eq:unscaled-target}
  \end{equation}
  holds for every graphon $W$ if and only if $a\le a_k^\star$.
\end{theorem}

The critical equation may also be written as
\begin{equation}
  (2a_k^\star)^{2k-1}(2a_k^\star+4k)=4k+2.
  \label{eq:critical-equivalent}
\end{equation}
For $k=2$, this becomes
\[
  8(a_2^\star)^3(a_2^\star+4)=5,
\]
and
\[
  a_2^\star=0.51721344027009436197\ldots.
\]
Thus Theorem~\ref{thm:main} includes the exact determination of
$\pi(C_4,C_5)$.

\section{Finite-rank reduction and the trace model}

We first reduce the proof to step graphons.  The details are included so that no
compactness theorem for graphons is needed.

\begin{lemma}[Step-graphon reduction]
  \label{lem:step-reduction}
  Fix integers $r_1,r_2\ge3$ and a continuous function
  $\Phi\colon\R^2\to\R$.  Suppose that
  \[
    \Phi\bigl(t(C_{r_1},W),t(C_{r_2},1-W)\bigr)\ge0
  \]
  for every symmetric step graphon $W$.  Then the same inequality holds for
  every graphon $W$.
\end{lemma}

\begin{proof}
  For $m\ge1$, partition $[0,1]$ into the $2^m$ dyadic intervals of length
  $2^{-m}$, and let $\mathcal A_m$ be the finite sigma-algebra generated by
  this partition.  Set
  \[
    W_m:=\mathbb E[W\mid \mathcal A_m\otimes\mathcal A_m].
  \]
  Conditional expectation preserves the bounds $0\le W\le1$, and symmetry is
  preserved because $\mathcal A_m\otimes\mathcal A_m$ is invariant under
  interchanging the two coordinates.  Thus $W_m$ is a symmetric step graphon.

  The spaces of dyadic rectangle step functions are nested and their union is
  dense in $L^2([0,1]^2)$.  Given $\varepsilon>0$, choose such a step function
  $G$, measurable with respect to $\mathcal A_M\otimes\mathcal A_M$, for which
  $\|W-G\|_2<\varepsilon$.  For $m\ge M$, the $L^2$-contraction property of
  conditional expectation gives
  \[
    \|W_m-G\|_2
    =\|\mathbb E[W-G\mid\mathcal A_m\otimes\mathcal A_m]\|_2
    \le\|W-G\|_2<\varepsilon.
  \]
  Hence $\|W_m-W\|_2<2\varepsilon$ for all $m\ge M$, so
  $\|W_m-W\|_2\to0$.  In particular, $\|W_m-W\|_1\to0$.  Also
  \[
    1-W_m
    =\mathbb E[1-W\mid\mathcal A_m\otimes\mathcal A_m],
  \]
  so $\|(1-W_m)-(1-W)\|_1\to0$.

  We claim that for every fixed integer $r\ge3$,
  \begin{equation}
    |t(C_r,W_m)-t(C_r,W)|
    \le r\|W_m-W\|_1.
    \label{eq:cycle-continuity}
  \end{equation}
  Indeed, write $x_{r+1}=x_1$ and use the telescoping identity
  \[
    \prod_{i=1}^r a_i-\prod_{i=1}^r b_i
    =
    \sum_{j=1}^r
    \left(\prod_{i<j}a_i\right)(a_j-b_j)
    \left(\prod_{i>j}b_i\right).
  \]
  Taking
  \[
    a_i=W_m(x_i,x_{i+1}),
    \qquad
    b_i=W(x_i,x_{i+1}),
  \]
  and using $0\le a_i,b_i\le1$, integration gives
  \eqref{eq:cycle-continuity}.  The same argument applies to $1-W_m$ and
  $1-W$.  Thus both cycle densities in the asserted inequality converge.
  Continuity of $\Phi$ now lets the step-graphon inequality pass to the limit.
\end{proof}

Now let $K$ be a symmetric step kernel, constant on every rectangle
$B_i\times B_j$ of a finite measurable partition
\[
  [0,1]=B_1\sqcup\cdots\sqcup B_N,
  \qquad |B_i|>0.
\]
Let $T_K$ be the associated integral operator,
\[
  (T_Kf)(x)=\int_0^1K(x,y)f(y)\,dy.
\]
Put
\[
  E:=\operatorname{span}\{\one_{B_1},\ldots,\one_{B_N}\}.
\]
Then $T_K$ maps $L^2([0,1])$ into $E$ and vanishes on $E^\perp$.  In the
orthonormal basis
\[
  e_i:=\frac{\one_{B_i}}{\sqrt{|B_i|}},
\]
the restriction $T_K|_E$ is represented by the real symmetric matrix
\[
  \bigl(K_{ij}\sqrt{|B_i||B_j|}\bigr)_{i,j=1}^N,
\]
where $K_{ij}$ is the value of $K$ on $B_i\times B_j$.

\begin{lemma}[Cycle densities are traces]
  \label{lem:trace-density}
  For every integer $r\ge3$,
  \begin{equation}
    t(C_r,K)=\tr(T_K^r),
    \label{eq:trace-density}
  \end{equation}
  where the trace is taken on the finite-dimensional space $E$.
\end{lemma}

\begin{proof}
  Expanding the matrix trace in the basis $(e_i)_{i=1}^N$ gives
  \[
    \tr(T_K^r)
    =
    \sum_{i_1,\ldots,i_r=1}^N
    |B_{i_1}|\cdots|B_{i_r}|
    K_{i_1i_2}K_{i_2i_3}\cdots K_{i_ri_1},
  \]
  which is exactly the step-function expansion of $t(C_r,K)$.
\end{proof}

Let $J:=T_1$ be the operator of the constant kernel $1$.  Since the underlying
measure space has total mass one,
\[
  Jf=\langle \one,f\rangle\one,
\]
so $J$ is the rank-one orthogonal projection onto $\R\one$.  If $U=1-W$ and
$T:=T_U$, then
\begin{equation}
  T_W=J-T.
  \label{eq:complement-operator}
\end{equation}

\section{A rank-one Lidskii--Karamata inequality}

For a vector $x\in\R^N$, write $x^\downarrow$ for its coordinates rearranged
in nonincreasing order.  Recall that $x$ is \emph{majorized} by $y$, written
$x\prec y$, if
\[
  \sum_{i=1}^j x_i^\downarrow
  \le
  \sum_{i=1}^j y_i^\downarrow
  \quad(1\le j<N),
  \qquad
  \sum_{i=1}^N x_i=\sum_{i=1}^N y_i.
\]
Karamata's inequality (see, e.g., \cite[Chapter~4]{MOA}) says that if
$x\prec y$ and $\varphi\colon\R\to\R$ is convex, then
\begin{equation}
  \sum_{i=1}^N\varphi(x_i)
  \le
  \sum_{i=1}^N\varphi(y_i).
  \label{eq:karamata}
\end{equation}

The next lemma is the special rank-one form of Lidskii's additive
majorization theorem that we need; for the general form see
\cite[Chapter~III]{Bhatia} or \cite[Chapter~9]{MOA}.  We include a
self-contained proof.

\begin{lemma}[Rank-one lower majorization]
  \label{lem:rank-one-majorization}
  Let $A$ be an $N\times N$ real symmetric matrix with eigenvalues
  \[
    \alpha_1\ge\alpha_2\ge\cdots\ge\alpha_N,
  \]
  and let $P$ be a rank-one orthogonal projection.  Let
  \[
    \mu_1\ge\mu_2\ge\cdots\ge\mu_N
  \]
  be the eigenvalues of $A+P$.  Then
  \begin{equation}
    (\alpha_1,\ldots,\alpha_{N-1},\alpha_N+1)
    \prec
    (\mu_1,\ldots,\mu_N).
    \label{eq:rank-one-majorization}
  \end{equation}
\end{lemma}

\begin{proof}
  The eigenvalues of a positive rank-one perturbation interlace:
  \begin{equation}
    \mu_i\ge\alpha_i\ge\mu_{i+1}
    \qquad(1\le i<N),
    \label{eq:rank-one-interlace}
  \end{equation}
  and also $\mu_N\ge\alpha_N$.  This follows directly from the
  Courant--Fischer min--max principle; equivalently, it is the rank-one case of
  Weyl interlacing (see, e.g., \cite[Chapter~III]{Bhatia}).

  Set
  \[
    v:=(\alpha_1,\ldots,\alpha_{N-1},\alpha_N+1).
  \]
  Fix $1\le j<N$.  The sum of the $j$ largest coordinates of $v$ is the
  maximum of $\sum_{i\in S}v_i$ over all $j$-element subsets
  $S\subseteq\{1,\ldots,N\}$.

  If $N\notin S$, then
  \[
    \sum_{i\in S}v_i
    \le\sum_{i=1}^j\alpha_i
    \le\sum_{i=1}^j\mu_i,
  \]
  by the first inequality in \eqref{eq:rank-one-interlace}.

  If $N\in S$, then
  \begin{equation}
    \sum_{i\in S}v_i
    \le
    1+\alpha_N+\sum_{i=1}^{j-1}\alpha_i.
    \label{eq:subset-with-N}
  \end{equation}
  Since $\tr(A+P)=\tr(A)+1$ and
  $\mu_{i+1}\le\alpha_i$ for $j\le i<N$, we have
  \begin{align*}
    \sum_{i=1}^j\mu_i
    &=
    \tr(A+P)-\sum_{i=j+1}^N\mu_i\\
    &\ge
    1+\sum_{i=1}^N\alpha_i-
    \sum_{i=j}^{N-1}\alpha_i\\
    &=
    1+\alpha_N+\sum_{i=1}^{j-1}\alpha_i.
  \end{align*}
  This dominates the right side of \eqref{eq:subset-with-N}.  Thus the first
  $j$ majorization inequality holds.  Equality of the total sums follows from
  $\tr(A+P)=\tr(A)+1$.
\end{proof}

\begin{corollary}[Even trace under subtraction of a rank-one projection]
  \label{cor:rank-one-trace}
  Let $T$ be a real symmetric $N\times N$ matrix with eigenvalues
  \[
    \lambda_1\ge\lambda_2\ge\cdots\ge\lambda_N,
  \]
  and let $P$ be a rank-one orthogonal projection.  If $n\ge2$ is even, then
  \begin{equation}
    \tr(P-T)^n
    \ge
    (1-\lambda_1)^n+\sum_{i=2}^N|\lambda_i|^n.
    \label{eq:rank-one-trace-bound}
  \end{equation}
\end{corollary}

\begin{proof}
  Apply Lemma~\ref{lem:rank-one-majorization} to $A=-T$.  The decreasingly
  ordered eigenvalues of $-T$ are
  \[
    -\lambda_N\ge\cdots\ge-\lambda_1.
  \]
  Hence
  \[
    (-\lambda_N,\ldots,-\lambda_2,1-\lambda_1)
    \prec
    \lambda(P-T),
  \]
  where $\lambda(P-T)$ denotes the eigenvalue vector of $P-T$.  Apply
  Karamata's inequality to the convex function $x\mapsto x^n$.  Since $n$ is
  even, $(-\lambda_i)^n=|\lambda_i|^n$, and
  \eqref{eq:rank-one-trace-bound} follows.
\end{proof}

\section{Spectral constraints for graphon operators}

The following elementary bounds use the pointwise condition $0\le U\le1$.
They do not assume that $T_U$ is positive semidefinite.

\begin{lemma}[Perron domination and the tail square budget]
  \label{lem:spectral-budget}
  Let $U$ be a symmetric step graphon, choose a finite partition on which $U$
  is constant, let $E$ be its associated step-function space, and let
  $T=T_U|_E$.  Let
  \[
    \lambda_1\ge\lambda_2\ge\cdots\ge\lambda_N
  \]
  be the eigenvalues of $T$.  Then
  \begin{align}
    &0\le\lambda_1\le1,
    \label{eq:lambda1-range}\\
    &|\lambda_i|\le\lambda_1
    \quad(1\le i\le N),
    \label{eq:perron-domination}\\
    &\sum_{i=2}^N\lambda_i^2
    \le\lambda_1(1-\lambda_1)
    \le\frac14.
    \label{eq:tail-square-budget}
  \end{align}
\end{lemma}

\begin{proof}
  Let $f$ be a real unit vector in $E$.  Since the elements of $E$ are
  constant on every cell of the partition, $|f|$ is again a unit vector in
  $E$.  Pointwise nonnegativity of $U$ gives
  \begin{equation}
    |\langle f,Tf\rangle|
    \le
    \langle |f|,T|f|\rangle,
    \label{eq:absolute-rayleigh}
  \end{equation}
  and by the Rayleigh principle the right side is at most $\lambda_1$.
  Applying \eqref{eq:absolute-rayleigh} to a unit eigenvector of $\lambda_i$
  shows that $|\lambda_i|\le\lambda_1$, proving
  \eqref{eq:perron-domination}.

  Next, the pointwise bound $U\le1$ gives
  \[
    \langle |f|,T|f|\rangle
    \le
    \left(\int_0^1|f|\right)^2,
  \]
  and the Cauchy--Schwarz inequality on the probability space $[0,1]$ gives
  $\int_0^1|f|\le\|f\|_2=1$.  Hence $\lambda_1\le1$.  Also $\lambda_1\ge0$,
  for example by testing the nonnegative unit vector $\one$.  This proves
  \eqref{eq:lambda1-range}.

  Finally, put
  \[
    s:=\int_{[0,1]^2}U(x,y)\,dx\,dy
    =\langle\one,T\one\rangle.
  \]
  Since $\|\one\|_2=1$, the Rayleigh principle gives $s\le\lambda_1$.  On the
  other hand,
  \[
    \sum_{i=1}^N\lambda_i^2
    =\tr(T^2)
    =\int_{[0,1]^2}U(x,y)^2\,dx\,dy
    \le s.
  \]
  Hence
  \[
    \sum_{i=2}^N\lambda_i^2
    \le s-\lambda_1^2
    \le\lambda_1-\lambda_1^2
    \le\frac14,
  \]
  as claimed.
\end{proof}

\section{The scalar branch and the critical parameter}

Fix an even integer
\[
  n\ge4.
\]
Eventually we set $n=2k$.  For $a\in(0,1)$, put
\[
  b:=1-a,
\]
and define
\begin{equation}
  \kappa=\kappa_n(a)
  :=\frac{n a^{n-1}}{(n+1)b^n},
  \qquad
  \rho_n(a)
  :=\frac{a^{n-1}(a+n)}{n+1}.
  \label{eq:kappa-rho}
\end{equation}
Multiplying \eqref{eq:unscaled-target}, with $n=2k$, by $na^{n-1}$ shows
that the desired commonality inequality is equivalent to
\begin{equation}
  t(C_n,W)+\kappa\,t(C_{n+1},1-W)
  \ge\rho_n(a).
  \label{eq:scaled-target}
\end{equation}
Indeed,
\[
  na^{n-1}
  \left(\frac an+\frac b{n+1}\right)
  =a^n+\frac{n}{n+1}a^{n-1}b
  =\frac{a^{n-1}(a+n)}{n+1}.
\]

Define
\begin{equation}
  f_\kappa(x):=(1-x)^n+\kappa x^{n+1},
  \qquad 0\le x\le1.
  \label{eq:f-kappa}
\end{equation}
Then
\begin{align*}
  f_\kappa'(x)
  &=-n(1-x)^{n-1}+\kappa(n+1)x^n,\\
  f_\kappa''(x)
  &=n(n-1)(1-x)^{n-2}+\kappa n(n+1)x^{n-1}>0.
\end{align*}
By the definition of $\kappa$,
\[
  f_\kappa'(b)
  =-na^{n-1}+\kappa(n+1)b^n
  =0.
\]
Therefore $f_\kappa$ is strictly convex and has its unique minimum at $b$.
Moreover,
\begin{equation}
  \min_{0\le x\le1}f_\kappa(x)
  =f_\kappa(b)
  =a^n+\kappa b^{n+1}
  =\rho_n(a).
  \label{eq:scalar-minimum}
\end{equation}

\begin{lemma}[The critical point]
  \label{lem:critical-point}
  There is a unique $\alpha_n^\star\in(1/2,1)$ such that
  \begin{equation}
    \rho_n(\alpha_n^\star)=2^{1-n}.
    \label{eq:critical-n}
  \end{equation}
\end{lemma}

\begin{proof}
  Differentiation gives
  \[
    \rho_n'(a)
    =\frac{n a^{n-2}(a+n-1)}{n+1}>0.
  \]
  Also
  \[
    \rho_n\left(\frac12\right)
    =2^{1-n}\frac{n+1/2}{n+1}
    <2^{1-n},
    \qquad
    \rho_n(1)=1>2^{1-n}.
  \]
  The intermediate value theorem and strict monotonicity give the result.
\end{proof}

The next quantitative bounds are the only estimates on the critical point
needed in the proof.

\begin{lemma}[Bounds on the normalized coefficient]
  \label{lem:kappa-bounds}
  Let $\alpha_n^\star$ be as in Lemma~\ref{lem:critical-point}, and put
  \[
    b_n^\star:=1-\alpha_n^\star,
    \qquad
    \kappa_n^\star:=\kappa_n(\alpha_n^\star).
  \]
  Then
  \begin{equation}
    \kappa_n^\star b_n^\star<1,
    \qquad
    \kappa_n^\star<\frac{27}{13}<2\sqrt2.
    \label{eq:kappa-star-bounds}
  \end{equation}
  Consequently, for every $0<a\le\alpha_n^\star$, with $b=1-a$ and
  $\kappa=\kappa_n(a)$,
  \begin{equation}
    \boxed{\kappa b<1,
    \qquad
    \kappa<2\sqrt2.}
    \label{eq:kappa-bounds-all-a}
  \end{equation}
\end{lemma}

\begin{proof}
  Write
  \[
    \alpha_n^\star=\frac{1+\delta}{2},
    \qquad
    b_n^\star=\frac{1-\delta}{2},
    \qquad
    \delta\in(0,1).
  \]
  Equation \eqref{eq:critical-n} is equivalent to
  \begin{equation}
    (1+\delta)^{n-1}(2n+1+\delta)=2n+2.
    \label{eq:delta-critical}
  \end{equation}
  Therefore
  \[
    (1+\delta)^{n-1}
    <\frac{2n+2}{2n+1}
    =1+\frac{1}{2n+1}.
  \]
  Bernoulli's inequality gives
  \[
    1+(n-1)\delta
    \le(1+\delta)^{n-1},
  \]
  and hence
  \begin{equation}
    0<\delta<\frac{1}{(n-1)(2n+1)}\le\frac1{27},
    \label{eq:delta-bound}
  \end{equation}
  because $n\ge4$.

  Using \eqref{eq:delta-critical}, we compute
  \begin{align}
    \kappa_n^\star b_n^\star
    &=
    \frac{n}{n+1}
    \left(\frac{1+\delta}{1-\delta}\right)^{n-1}
    \notag\\
    &=
    \frac{2n}{(2n+1+\delta)(1-\delta)^{n-1}}.
    \label{eq:kappa-b-delta}
  \end{align}
  Again by Bernoulli's inequality and \eqref{eq:delta-bound},
  \[
    (1-\delta)^{n-1}
    \ge1-(n-1)\delta
    >1-\frac1{2n+1}
    =\frac{2n}{2n+1}.
  \]
  Thus the denominator in \eqref{eq:kappa-b-delta} is strictly larger than
  $2n$, proving $\kappa_n^\star b_n^\star<1$.

  It follows that
  \[
    \frac1{\kappa_n^\star}>b_n^\star
    =\frac{1-\delta}{2}
    >\frac{13}{27},
  \]
  where the last inequality uses \eqref{eq:delta-bound}.  Hence
  \[
    \kappa_n^\star<\frac{27}{13}<2\sqrt2.
  \]

  Finally,
  \[
    \kappa_n(a)
    =\frac{n}{n+1}\frac{a^{n-1}}{(1-a)^n}
  \]
  is strictly increasing in $a$, and so is
  \[
    \kappa_n(a)(1-a)
    =\frac{n}{n+1}
    \left(\frac{a}{1-a}\right)^{n-1}.
  \]
  Therefore the endpoint bounds remain valid for every $a\le\alpha_n^\star$.
\end{proof}

\section{Proof of the lower bound}

We now prove that \eqref{eq:scaled-target} holds whenever
$a\le\alpha_n^\star$.

\begin{proposition}[Spectral reduction]
  \label{prop:spectral-reduction}
  Let $n\ge4$ be even, let $\kappa>0$, and let $W$ be a symmetric step
  graphon.  Choose a finite partition on which $W$ is constant, let $E$ be its
  associated step-function space, put $U=1-W$, and let $T=T_U|_E$.  Order the
  eigenvalues of $T$ as
  \[
    \lambda_1\ge\lambda_2\ge\cdots\ge\lambda_N.
  \]
  Then
  \begin{equation}
    t(C_n,W)+\kappa t(C_{n+1},U)
    \ge
    f_\kappa(\lambda_1)
    +
    \sum_{i=2}^N\lambda_i^n(1+\kappa\lambda_i).
    \label{eq:master-spectral}
  \end{equation}
\end{proposition}

\begin{proof}
  Let $P:=J|_E$.  Since $\one\in E$ and $\|\one\|_2=1$, the operator $P$ is a
  rank-one orthogonal projection on $E$.  Moreover,
  $T_W|_E=P-T$ by \eqref{eq:complement-operator}.  Therefore, by
  Lemma~\ref{lem:trace-density} and Corollary~\ref{cor:rank-one-trace},
  \[
    t(C_n,W)
    =\tr(P-T)^n
    \ge
    (1-\lambda_1)^n+
    \sum_{i=2}^N|\lambda_i|^n.
  \]
  Since $n$ is even, $|\lambda_i|^n=\lambda_i^n$.  Also
  \[
    t(C_{n+1},U)=\tr(T^{n+1})=\sum_{i=1}^N\lambda_i^{n+1}.
  \]
  Adding the two expressions gives
  \begin{align*}
    t(C_n,W)+\kappa t(C_{n+1},U)
    &\ge
    (1-\lambda_1)^n+\kappa\lambda_1^{n+1}\\
    &\quad+
    \sum_{i=2}^N
    \bigl(\lambda_i^n+\kappa\lambda_i^{n+1}\bigr),
  \end{align*}
  which is exactly \eqref{eq:master-spectral}.
\end{proof}

\begin{proposition}[Commonality up to the critical point]
  \label{prop:lower-bound}
  Let $n\ge4$ be even and $0<a\le\alpha_n^\star$.  Put $b=1-a$ and let
  $\kappa$ be as in \eqref{eq:kappa-rho}.  Then every graphon $W$ satisfies
  \begin{equation}
    t(C_n,W)+\kappa t(C_{n+1},1-W)
    \ge\rho_n(a).
    \label{eq:lower-bound-final}
  \end{equation}
\end{proposition}

\begin{proof}
  Apply Lemma~\ref{lem:step-reduction} with
  $\Phi(x,y)=x+\kappa y-\rho_n(a)$.  It is enough to prove the assertion for a
  symmetric step graphon.  Choose a finite partition on which $W$ is constant,
  let $E$ be its associated step-function space, set $U=1-W$, and let
  $T=T_U|_E$.  Write its eigenvalues as
  \[
    \lambda_1\ge\lambda_2\ge\cdots\ge\lambda_N.
  \]
  By Proposition~\ref{prop:spectral-reduction},
  \begin{equation}
    t(C_n,W)+\kappa t(C_{n+1},U)
    \ge
    f_\kappa(\lambda_1)
    +
    \sum_{i=2}^N\lambda_i^n(1+\kappa\lambda_i).
    \label{eq:lower-start}
  \end{equation}

  Call an eigenvalue $\lambda_i$ with $i\ge2$ \emph{dangerous} if
  \begin{equation}
    1+\kappa\lambda_i<0,
    \qquad\text{equivalently}\qquad
    \lambda_i<-\frac1\kappa.
    \label{eq:dangerous-definition}
  \end{equation}
  Because $n$ is even, every nondangerous summand
  $\lambda_i^n(1+\kappa\lambda_i)$ in \eqref{eq:lower-start} is nonnegative.

  \medskip
  \noindent\emph{Case 1: no dangerous eigenvalue exists.}
  Then \eqref{eq:lower-start} and \eqref{eq:scalar-minimum} give
  \[
    t(C_n,W)+\kappa t(C_{n+1},U)
    \ge f_\kappa(\lambda_1)
    \ge f_\kappa(b)
    =\rho_n(a).
  \]

  \medskip
  \noindent\emph{Case 2: a dangerous eigenvalue exists.}
  There is at most one dangerous eigenvalue.  Indeed, if two eigenvalues were
  smaller than $-1/\kappa$, then Lemma~\ref{lem:spectral-budget} and
  Lemma~\ref{lem:kappa-bounds} would give
  \[
    \frac14
    \ge\sum_{i=2}^N\lambda_i^2
    >\frac{2}{\kappa^2}
    >\frac14,
  \]
  a contradiction.  Write the unique dangerous eigenvalue as
  \[
    -\beta,
    \qquad
    \beta>\frac1\kappa.
  \]
  All remaining terms in the sum in \eqref{eq:lower-start} are nonnegative.
  Perron domination \eqref{eq:perron-domination} gives
  \begin{equation}
    \lambda_1\ge\beta.
    \label{eq:perron-pays}
  \end{equation}
  Moreover, Lemma~\ref{lem:kappa-bounds} gives
  \begin{equation}
    \beta>\frac1\kappa>b.
    \label{eq:beta-above-b}
  \end{equation}
  By \eqref{eq:perron-pays} and \eqref{eq:lambda1-range},
  \[
    0<\beta\le\lambda_1\le1.
  \]
  Since $f_\kappa$ is strictly convex and minimized at $b$, it is increasing
  on $[b,1]$.  Hence, by \eqref{eq:perron-pays} and
  \eqref{eq:beta-above-b},
  \[
    f_\kappa(\lambda_1)\ge f_\kappa(\beta).
  \]
  Therefore \eqref{eq:lower-start} yields
  \begin{align*}
    t(C_n,W)+\kappa t(C_{n+1},U)
    &\ge
    f_\kappa(\beta)+\beta^n(1-\kappa\beta)\\
    &=
    (1-\beta)^n+\kappa\beta^{n+1}
    +\beta^n-\kappa\beta^{n+1}\\
    &=
    (1-\beta)^n+\beta^n.
  \end{align*}
  Since $n$ is even, convexity of $x\mapsto x^n$ gives
  \begin{equation}
    (1-\beta)^n+\beta^n
    \ge2\left(\frac12\right)^n
    =2^{1-n}.
    \label{eq:two-point-convexity}
  \end{equation}
  Finally, $a\le\alpha_n^\star$ and the monotonicity of $\rho_n$ imply
  \[
    \rho_n(a)\le\rho_n(\alpha_n^\star)=2^{1-n}.
  \]
  Combining this with \eqref{eq:two-point-convexity} proves
  \eqref{eq:lower-bound-final}.
\end{proof}

\section{The balanced two-clique obstruction and completion}

Let $W_{\mathrm{cl}}$ be the balanced disjoint union of two cliques:
\[
  W_{\mathrm{cl}}(x,y)
  :=
  \begin{cases}
    1,
    &x,y\in[0,1/2]
      \text{ or }x,y\in(1/2,1],\\
    0,&\text{otherwise}.
  \end{cases}
\]
Since $C_n$ is connected, every homomorphism contributing to
$t(C_n,W_{\mathrm{cl}})$ maps all of its vertices into the same half.  Hence
\begin{equation}
  t(C_n,W_{\mathrm{cl}})
  =2\left(\frac12\right)^n
  =2^{1-n}.
  \label{eq:two-clique-even-cycle}
\end{equation}
The complement $1-W_{\mathrm{cl}}$ is the balanced complete bipartite graphon.
Since $n+1$ is odd,
\begin{equation}
  t(C_{n+1},1-W_{\mathrm{cl}})=0.
  \label{eq:bipartite-odd-cycle}
\end{equation}
Thus
\begin{equation}
  t(C_n,W_{\mathrm{cl}})
  +\kappa t(C_{n+1},1-W_{\mathrm{cl}})
  =2^{1-n}.
  \label{eq:obstruction-value}
\end{equation}

If $a>\alpha_n^\star$, then strict monotonicity of $\rho_n$ gives
\[
  \rho_n(a)>\rho_n(\alpha_n^\star)=2^{1-n}.
\]
Therefore $W_{\mathrm{cl}}$ violates the scaled commonality inequality
\eqref{eq:scaled-target}.  It follows that
\[
  \pi(C_n,C_{n+1})\subseteq(0,\alpha_n^\star].
\]
Proposition~\ref{prop:lower-bound} gives the reverse inclusion.  Taking
$n=2k$, the number $\alpha_{2k}^\star$ is exactly the number
$a_k^\star$ defined in Theorem~\ref{thm:main}, so the theorem follows.

\begin{remark}[The endpoint transition]
  For every $a\in(0,1)$, the constant graphon $W\equiv a$ attains equality in
  the defining commonality inequality.  At $a=\alpha_n^\star$, the balanced
  two-clique graphon also attains equality because
  $\rho_n(\alpha_n^\star)=2^{1-n}$.  Thus the endpoint is exactly the point where
  the constant branch meets the balanced two-clique branch.
\end{remark}

\begin{remark}[The role of the rank-one decomposition]
  The identity $T_W=J-T_{1-W}$ separates the one-dimensional constant state
  from the remaining spectrum.  Corollary~\ref{cor:rank-one-trace} is a global
  form of this separation: it accounts for every coupling between the
  constant state and the orthogonal directions at once.  The tail square
  budget then leaves at most one negative mode for which
  $1+\kappa\lambda<0$, and the Perron eigenvalue compensates for that mode.
  For the balanced two-clique graphon the spectrum of $T_{1-W}|_E$ is
  $\{\tfrac12,-\tfrac12\}$ and \eqref{eq:rank-one-trace-bound} holds with
  equality, so Corollary~\ref{cor:rank-one-trace} is tight exactly at the
  extremal point.
\end{remark}

\begin{thebibliography}{9}

\bibitem{BMN}
N. Behague, N. Morrison, and J. A. Noel,
\emph{Common pairs of graphs},
Combin. Probab. Comput. \textbf{34} (2025), no.~5, 649--670.

\bibitem{Bhatia}
R. Bhatia,
\emph{Matrix Analysis},
Graduate Texts in Mathematics, vol.~169,
Springer, New York, 1997.

\bibitem{Lovasz}
L. Lov\'asz,
\emph{Large Networks and Graph Limits},
American Mathematical Society Colloquium Publications, vol.~60,
American Mathematical Society, Providence, RI, 2012.

\bibitem{MOA}
A. W. Marshall, I. Olkin, and B. C. Arnold,
\emph{Inequalities: Theory of Majorization and Its Applications},
second edition,
Springer Series in Statistics,
Springer, New York, 2011.

\end{thebibliography}

\end{document}
